\documentclass[11pt]{article}
\usepackage{definitions}

\begin{document}

\centerline{\textbf{\textsc{Reading the Coronary Microcirculation from How the Heart Muscle Shrinks with Each Beat}}}

\textbf{SPECIFIC AIMS} -- The heart muscle is fed by its own microcirculation, vessels inside the wall too small for any angiogram to show. When these vessels are diseased, patients have angina with open arteries: nearly two thirds of patients catheterized for chest pain have no blocked artery, and in many of them the small vessels are the cause; three quarters of patients with heart failure with preserved ejection fraction have diseased microvessels.\cite{kunadian20,delbuono21,patel10,shah18} Today's tests, by catheter or \gls*{pet}, measure only how much blood flow can rise under a drug mimicking exercise (the flow reserve).\cite{kunadian20,vink21} Each contraction squeezes nearly half the blood out of these vessels, and the wall shrinks with every beat, by 2 to 25\% of its volume depending on the method.\cite{judd91,ashikaga08,brown25} Whether chamber pressure or the muscle fibers around the vessels push the blood out, and how much of the shrinkage is blood, is unsettled after fifty years:\cite{westerhof06} no experiment has measured a region's volume, inflow, outflow, and pressures together. \textbf{If the shrinking of the muscle with each beat depends on how much blood the microvessels hold and how easily they are squeezed, every echocardiogram and \gls*{mri} already records a sign of microvascular disease that no one can yet read.}

Our \textbf{long-term goal} is to diagnose microvascular disease without a catheter, from how the heart muscle moves on routine images. Our \textbf{central hypothesis} is that the shrinkage is set by how hard the muscle contracts and by how much blood the microvessels hold and how easily they are squeezed, so it falls when the heart is weak or the microvessels diseased. In seven dogs we sewed tiny ultrasound transducers (crystals) into the muscle to track the volume of 1-cc cubes in regions fed by an artery we narrowed and by a healthy one, and recorded flow in the artery, blood reaching the muscle, and pressures in chamber and muscle.\cite{stendahl23} The narrowing shrank the relaxed volume of the starved cube by 5\%, and of the healthy cube as much although its blood supply was unchanged: a change the flow measurement missed, which we think reflects the microvessels. Pfaller (PI) writes computer simulations of blood flow and the heart wall tuned to data;\cite{pfaller24,richter25} Stendahl (PI) built the dog experiments with co-investigators Sinusas and Liu (blood-volume imaging)\cite{mohyuddin18,yousefi23} and Duncan (3D echocardiography).\cite{ta24} Stendahl's K08 award supports a machine that keeps human hearts turned down for transplant beating.

\aim{1}{To determine what sets the shrinking of the muscle with each beat.}
{We hypothesize that the muscle fibers tightening around the vessels, not the chamber pressure pushing through the wall, squeeze out most of the blood, and that the tissue itself compresses little. We will simulate the wall as nearly incompressible tissue holding arterial, capillary, and venous blood in vessels that swell with pressure,\cite{cookson12,lee16b} tune it to our dog data, test it on data it was not tuned to, and apportion the shrinkage among blood pushed back into the arteries, blood pushed into the veins, and tissue compression. Completing this aim will show which mechanism dominates and let heart models include the blood they now ignore.\cite{brown25,liu21c}}

\aim{2}{To measure how much blood each beat squeezes from the microvessels and what controls it.}
{We hypothesize that dilating the microvessels increases the shrinkage, weakening contraction decreases it, and blocking the microvessels abolishes it even with the artery open. In \note{XX} dogs of both sexes we will record the volume, inflow, and outflow of the region fed by an artery we narrow and the pressure inside the muscle, while dilating the microvessels (adenosine), weakening contraction (esmolol), strengthening it (dobutamine), or blocking the microvessels with tiny beads (embolization). On the machine we will vary blood supply and contraction of beating human hearts one at a time, then stop them to measure how vessels and tissue alone respond to pressure. Completing this aim will show how much of the shrinkage is blood leaving the vessels.}

\aim{3}{To validate noninvasive measurement of the shrinkage and test it in patients.}
{We hypothesize that three-dimensional echocardiography, checked against the crystals in the Aim 2 dogs, measures the shrinkage within \note{XX}\%, and that in patients with angina and open arteries its change under adenosine, which dilates the microvessels but barely changes contraction, identifies microvascular disease better than strain, the shortening it reports. We will then enroll \note{XX} of the \note{XX} women and men a year whose flow reserve Yale measures for angina with open arteries, and image each by echocardiography and \gls*{mri} at rest and under adenosine. Completing this aim will show whether routine examinations detect microvascular disease (flow reserve below 2, the consensus threshold)\cite{kunadian20} with accuracy (area under the curve) above \note{XX}.}

We will deliver the first computer model of heart muscle that holds blood, an accounting of where that blood goes with each beat, and a validated noninvasive measure of the shrinkage. \textbf{Patients with angina and no blockage on their angiogram could have their microvessels examined on the echocardiogram or \gls*{mri} they already get, and be treated for the cause rather than catheterized.}

\clearpage
\section*{\normalsize{REFERENCES CITED}}
\bibliographystyle{mrp}
\bibliography{martin_highlighted,new_references}
\end{document}
